% Основной текст отчёта.
% Графики из code/lab1.ipynb сохраняйте в папку tex/images
% под именами, указанными в плейсхолдерах ниже.

\newcommand{\imageplaceholder}[3]{%
  \begin{figure}[htbp]
    \centering
    \IfFileExists{images/#1}{%
      \includegraphics[width=0.9\textwidth]{#1}%
    }{%
      \fbox{%
        \parbox[c][5cm][c]{0.78\textwidth}{%
          \centering
          Место для рисунка\\[0.4em]
          \texttt{images/\detokenize{#1}}
        }%
      }%
    }
    \caption{#2}
    \label{#3}
  \end{figure}
}

\tableofcontents
\newpage

\section{Цель работы}

Исследовать положения равновесия нелинейных динамических систем, выполнить
их линеаризацию, определить тип и устойчивость изолированных положений
равновесия, построить фазовые портреты и синтезировать локально
стабилизирующие регуляторы.

\section{Исследование автономных нелинейных систем}

\subsection{Система 1}

\begin{equation}
  \begin{cases}
    \dot x_1=-x_1+2x_1^3+x_2,\\
    \dot x_2=-x_1-x_2
  \end{cases}
  \label{eq:system-1}
\end{equation}

Точки равновесия находятся из уравнений
\begin{equation}
  \begin{cases}
    -x_1+2x_1^3+x_2=0,\\
    -x_1-x_2=0
  \end{cases}
\end{equation}
Из второго уравнения $x_2=-x_1$. После подстановки в первое уравнение
получим
\begin{equation}
  2x_1^3-2x_1=2x_1(x_1-1)(x_1+1)=0
\end{equation}
Следовательно, система имеет три точки равновесия:
\begin{equation}
  \bm x^{(1)}=(0,0),\qquad
  \bm x^{(2)}=(1,-1),\qquad
  \bm x^{(3)}=(-1,1)
\end{equation}

Матрица Якоби системы имеет вид
\begin{equation}
  A_1(x_1,x_2)=
  \begin{pmatrix}
    -1+6x_1^2 & 1\\
    -1 & -1
  \end{pmatrix}
\end{equation}

Для малых отклонений
$\Delta\bm x=\bm x-\bm x^{(i)}$ линеаризованная система записывается как
\begin{equation}
  \Delta\dot{\bm x}=A_1\bigl(\bm x^{(i)}\bigr)\Delta\bm x
\end{equation}
В точке $\bm x^{(1)}=(0,0)$ получаем
\begin{equation}
  A_1(0,0)=
  \begin{pmatrix}
    -1 & 1\\
    -1 & -1
  \end{pmatrix},
  \qquad
  \lambda_{1,2}=-1\pm i
\end{equation}
Поэтому точка $(0,0)$ является устойчивым фокусом.

Для двух остальных точек матрицы линеаризации совпадают:
\begin{equation}
  A_1(1,-1)=A_1(-1,1)=
  \begin{pmatrix}
    5 & 1\\
    -1 & -1
  \end{pmatrix}
\end{equation}
Характеристическое уравнение и его корни равны
\begin{equation}
  \lambda^2-4\lambda-4=0,
  \qquad
  \lambda_{1,2}=2\pm2\sqrt{2}
\end{equation}
Одно собственное значение положительно, а другое отрицательно, поэтому
точки $(1,-1)$ и $(-1,1)$ являются неустойчивыми седлами.

\imageplaceholder{phase_portrait_1.png}{Фазовый портрет системы 1}{fig:phase-system-1}

\subsection{Система 2}

\begin{equation}
  \begin{cases}
    \dot x_1=x_1+x_1x_2,\\
    \dot x_2=-x_2+x_2^2+x_1x_2-x_1^3
  \end{cases}
  \label{eq:system-2}
\end{equation}

Точки равновесия находятся из уравнений
\begin{equation}
  \begin{cases}
    x_1+x_1x_2=0,\\
    -x_2+x_2^2+x_1x_2-x_1^3=0
  \end{cases}
\end{equation}
Матрица Якоби:
\begin{equation}
  A_2(x_1,x_2)=
  \begin{pmatrix}
    1+x_2 & x_1\\
    x_2-3x_1^2 & -1+2x_2+x_1
  \end{pmatrix}
\end{equation}

Из первого уравнения следует, что $x_1=0$ или $x_2=-1$. В результате
получаем три точки равновесия:
\begin{equation}
  \bm x^{(1)}=(0,0),\qquad
  \bm x^{(2)}=(0,1),\qquad
  \bm x^{(3)}=(1,-1)
\end{equation}
В точке $(0,0)$ имеем
\begin{equation}
  A_2(0,0)=
  \begin{pmatrix}
    1 & 0\\
    0 & -1
  \end{pmatrix},
  \qquad
  \lambda_1=1,\quad\lambda_2=-1
\end{equation}
Следовательно, точка $(0,0)$ является неустойчивым седлом. Для точки
$(0,1)$ получаем
\begin{equation}
  A_2(0,1)=
  \begin{pmatrix}
    2 & 0\\
    1 & 1
  \end{pmatrix},
  \qquad
  \lambda_1=2,\quad\lambda_2=1
\end{equation}
Оба собственных значения положительны, поэтому $(0,1)$ является
неустойчивым узлом. В точке $(1,-1)$
\begin{equation}
  \begin{aligned}
    A_2(1,-1)&=
    \begin{pmatrix}
      0 & 1\\
      -4 & -2
    \end{pmatrix}\\
    \det\bigl(A_2-\lambda I\bigr)
    &=
    \begin{vmatrix}
      -\lambda & 1\\
      -4 & -2-\lambda
    \end{vmatrix}
    =\lambda(\lambda+2)+4
    =\lambda^2+2\lambda+4=0\\
    D&=2^2-4\cdot 1\cdot 4=-12\\
    \lambda_{1,2}
    &=\frac{-2\pm\sqrt{-12}}{2}
    =-1\pm i\sqrt{3}
  \end{aligned}
\end{equation}
Таким образом, точка $(1,-1)$ является устойчивым фокусом.

\imageplaceholder{phase_portrait_2.png}{Фазовый портрет системы 2}{fig:phase-system-2}

\subsection{Система 3}

\begin{equation}
  \begin{cases}
    \dot x_1=x_2,\\
    \dot x_2=-x_1+x_2\left(1-x_1^2+0{,}1x_1^4\right)
  \end{cases}
  \label{eq:system-3}
\end{equation}

Точки равновесия находятся из уравнений
\begin{equation}
  \begin{cases}
    x_2=0,\\
    -x_1+x_2\left(1-x_1^2+0{,}1x_1^4\right)=0
  \end{cases}
\end{equation}
Матрица Якоби:
\begin{equation}
  A_3(x_1,x_2)=
  \begin{pmatrix}
    0 & 1\\
    -1+x_2(-2x_1+0{,}4x_1^3) & 1-x_1^2+0{,}1x_1^4
  \end{pmatrix}
\end{equation}

Из первого уравнения $x_2=0$, после чего второе уравнение даёт $x_1=0$.
Единственная точка равновесия равна
\begin{equation}
  \bm x^{(1)}=(0,0)
\end{equation}
Матрица линеаризации и её собственные значения:
\begin{equation}
  \begin{aligned}
    A_3(0,0)&=
    \begin{pmatrix}
      0 & 1\\
      -1 & 1
    \end{pmatrix}\\
    \det\bigl(A_3-\lambda I\bigr)
    &=
    \begin{vmatrix}
      -\lambda & 1\\
      -1 & 1-\lambda
    \end{vmatrix}
    =-\lambda(1-\lambda)+1
    =\lambda^2-\lambda+1=0\\
    D&=(-1)^2-4\cdot 1\cdot 1=-3\\
    \lambda_{1,2}
    &=\frac{1\pm\sqrt{-3}}{2}
    =\frac{1\pm i\sqrt{3}}{2}
  \end{aligned}
\end{equation}
Действительные части собственных значений положительны, поэтому точка
$(0,0)$ является неустойчивым фокусом.

\imageplaceholder{phase_portrait_3.png}{Фазовый портрет системы 3}{fig:phase-system-3}

\subsection{Система 4}

\begin{equation}
  \begin{cases}
    \dot x_1=(x_1-x_2)(1-x_1^2-x_2^2),\\
    \dot x_2=(x_1+x_2)(1-x_1^2-x_2^2)
  \end{cases}
  \label{eq:system-4}
\end{equation}

Точки равновесия находятся из уравнений
\begin{equation}
  \begin{cases}
    (x_1-x_2)(1-x_1^2-x_2^2)=0,\\
    (x_1+x_2)(1-x_1^2-x_2^2)=0
  \end{cases}
\end{equation}
Условия выполняются либо в начале координат, либо при
$x_1^2+x_2^2=1$. Поэтому множество равновесий имеет вид
\begin{equation}
  \bm x^{(1)}=(0,0),\qquad
  \mathcal E=\{(x_1,x_2):x_1^2+x_2^2=1\}
\end{equation}
Обозначив $q=1-x_1^2-x_2^2$, получим матрицу Якоби
\begin{equation}
  A_4(x_1,x_2)=
  \begin{pmatrix}
    q-2x_1(x_1-x_2) & -q-2x_2(x_1-x_2)\\
    q-2x_1(x_1+x_2) & q-2x_2(x_1+x_2)
  \end{pmatrix}
\end{equation}

Для дополнительного исследования поведения системы выполним переход
к полярным координатам

\begin{equation}
  x_1=r\cos\varphi,\qquad x_2=r\sin\varphi,
\end{equation}

где

\begin{equation}
  r^2=x_1^2+x_2^2.
\end{equation}

Дифференцируя последнее равенство по времени, получаем

\begin{equation}
  2r\dot r=2x_1\dot x_1+2x_2\dot x_2,
\end{equation}

откуда при $r>0$

\begin{equation}
  \dot r=\frac{x_1\dot x_1+x_2\dot x_2}{r}.
\end{equation}

Для угловой координаты продифференцируем соотношения

\[
x_1=r\cos\varphi,\qquad x_2=r\sin\varphi.
\]

Имеем

\begin{equation}
  \begin{aligned}
    \dot x_1&=\dot r\cos\varphi-r\dot\varphi\sin\varphi,\\
    \dot x_2&=\dot r\sin\varphi+r\dot\varphi\cos\varphi.
  \end{aligned}
\end{equation}

Рассмотрим комбинацию $x_1\dot x_2-x_2\dot x_1$. После подстановки
выражений для $x_1$, $x_2$, $\dot x_1$ и $\dot x_2$ получаем

\begin{equation}
  x_1\dot x_2-x_2\dot x_1
  =r^2\dot\varphi,
\end{equation}

следовательно,

\begin{equation}
  \dot\varphi=
  \frac{x_1\dot x_2-x_2\dot x_1}{r^2}.
\end{equation}

Таким образом, при $r>0$

\begin{equation}
  \dot r=\frac{x_1\dot x_1+x_2\dot x_2}{r},\qquad
  \dot\varphi=\frac{x_1\dot x_2-x_2\dot x_1}{r^2}.
\end{equation}

Поскольку

\[
x_1^2+x_2^2=r^2,
\]

ранее введенная величина $q$ в полярных координатах принимает вид

\begin{equation}
  q=1-r^2.
\end{equation}

Подставим систему~\eqref{eq:system-4} в выражение для $\dot r$:

\begin{equation}
  \begin{aligned}
    \dot r
    &=
    \frac{
      x_1(x_1-x_2)q+
      x_2(x_1+x_2)q
    }{r}\\
    &=
    \frac{
      \left(
        x_1^2-x_1x_2+x_1x_2+x_2^2
      \right)q
    }{r}\\
    &=
    \frac{x_1^2+x_2^2}{r}q
    =rq.
  \end{aligned}
\end{equation}

Следовательно,

\begin{equation}
  \dot r=r(1-r^2).
\end{equation}

Аналогично для угловой координаты

\begin{equation}
  \begin{aligned}
    \dot\varphi
    &=
    \frac{
      x_1(x_1+x_2)q-
      x_2(x_1-x_2)q
    }{r^2}\\
    &=
    \frac{
      \left(
        x_1^2+x_1x_2-x_1x_2+x_2^2
      \right)q
    }{r^2}\\
    &=
    \frac{x_1^2+x_2^2}{r^2}q
    =q.
  \end{aligned}
\end{equation}

Таким образом, исходная система при $r>0$ в полярных координатах
принимает вид

\begin{equation}
  \boxed{
  \begin{aligned}
    \dot r&=r(1-r^2),\\
    \dot\varphi&=1-r^2.
  \end{aligned}}
\end{equation}

Полученное представление позволяет непосредственно исследовать
радиальное движение траекторий. При $0<r<1$ выполняется
$\dot r>0$, поэтому траектории удаляются от начала координат.
При $r>1$ имеем $\dot r<0$, поэтому траектории движутся в сторону
окружности $r=1$. На самой окружности

\[
\dot r=0,\qquad \dot\varphi=0,
\]

то есть каждая ее точка является положением равновесия.

Заметим, что формула для $\dot\varphi$ содержит $r^2$ в знаменателе
и поэтому непосредственно неприменима при $r=0$. Устойчивость
начала координат исследуется отдельно с помощью линеаризации.

В начале координат
\begin{equation}
  \begin{aligned}
    A_4(0,0)&=
    \begin{pmatrix}
      1 & -1\\
      1 & 1
    \end{pmatrix}\\
    \det\bigl(A_4-\lambda I\bigr)
    &=
    \begin{vmatrix}
      1-\lambda & -1\\
      1 & 1-\lambda
    \end{vmatrix}
    =(1-\lambda)^2+1=0\\
    (1-\lambda)^2&=-1\\
    \lambda_{1,2}&=1\pm i
  \end{aligned}
\end{equation}
Следовательно, начало координат является неустойчивым фокусом. При
$0<r<1$ выполняется $\dot r>0$, а при $r>1$ выполняется $\dot r<0$,
поэтому окружность $r=1$ притягивает соседние траектории. Данная окружность является множеством неизолированных устойчивых точек равновесия.

\imageplaceholder{phase_portrait_4.png}{Фазовый портрет системы 4}{fig:phase-system-4}

\subsection{Система 5}

\begin{equation}
  \begin{cases}
    \dot x_1=-x_1^3+x_2,\\
    \dot x_2=x_1-x_2^3
  \end{cases}
  \label{eq:system-5}
\end{equation}

Точки равновесия находятся из уравнений
\begin{equation}
  \begin{cases}
    -x_1^3+x_2=0,\\
    x_1-x_2^3=0
  \end{cases}
\end{equation}
Матрица Якоби:
\begin{equation}
  A_5(x_1,x_2)=
  \begin{pmatrix}
    -3x_1^2 & 1\\
    1 & -3x_2^2
  \end{pmatrix}
\end{equation}

Из условий равновесия следует $x_2=x_1^3$ и $x_1=x_2^3$. Отсюда
получаем три точки:
\begin{equation}
  \bm x^{(1)}=(0,0),\qquad
  \bm x^{(2)}=(1,1),\qquad
  \bm x^{(3)}=(-1,-1)
\end{equation}
В начале координат
\begin{equation}
  A_5(0,0)=
  \begin{pmatrix}
    0 & 1\\
    1 & 0
  \end{pmatrix},
  \qquad
  \lambda_1=1,\quad\lambda_2=-1
\end{equation}
Точка $(0,0)$ является неустойчивым седлом. Для двух остальных точек
матрицы линеаризации совпадают:
\begin{equation}
  \begin{aligned}
    A_5(1,1)=A_5(-1,-1)&=
    \begin{pmatrix}
      -3 & 1\\
      1 & -3
    \end{pmatrix}\\
    \det\bigl(A_5-\lambda I\bigr)
    &=
    \begin{vmatrix}
      -3-\lambda & 1\\
      1 & -3-\lambda
    \end{vmatrix}
    =(-3-\lambda)^2-1=0\\
    (\lambda+3)^2-1
    &=(\lambda+2)(\lambda+4)=0\\
    \lambda_1&=-2,\qquad \lambda_2=-4
  \end{aligned}
\end{equation}
Следовательно, точки $(1,1)$ и $(-1,-1)$ являются устойчивыми узлами.

\imageplaceholder{phase_portrait_5.png}{Фазовый портрет системы 5}{fig:phase-system-5}

\subsection{Система 6}

\begin{equation}
  \begin{cases}
    \dot x_1=-x_1^3+x_2^3,\\
    \dot x_2=x_2^3x_1-x_2^3
  \end{cases}
  \label{eq:system-6}
\end{equation}

Точки равновесия находятся из уравнений
\begin{equation}
  \begin{cases}
    -x_1^3+x_2^3=0,\\
    x_2^3(x_1-1)=0
  \end{cases}
\end{equation}
Матрица Якоби:
\begin{equation}
  A_6(x_1,x_2)=
  \begin{pmatrix}
    -3x_1^2 & 3x_2^2\\
    x_2^3 & 3x_2^2(x_1-1)
  \end{pmatrix}
\end{equation}

Из первого уравнения следует $x_2=x_1$. После подстановки во второе
уравнение получаем две точки равновесия:
\begin{equation}
  \bm x^{(1)}=(0,0),\qquad
  \bm x^{(2)}=(1,1)
\end{equation}
В начале координат
\begin{equation}
  A_6(0,0)=
  \begin{pmatrix}
    0 & 0\\
    0 & 0
  \end{pmatrix},
  \qquad
  \lambda_1=\lambda_2=0
\end{equation}
Поэтому метод линеаризации не позволяет определить устойчивость точки
$(0,0)$. Для точки $(1,1)$
\begin{equation}
  \begin{aligned}
    A_6(1,1)&=
    \begin{pmatrix}
      -3 & 3\\
      1 & 0
    \end{pmatrix}\\
    \det\bigl(A_6-\lambda I\bigr)
    &=
    \begin{vmatrix}
      -3-\lambda & 3\\
      1 & -\lambda
    \end{vmatrix}
    =\lambda(\lambda+3)-3
    =\lambda^2+3\lambda-3=0\\
    D&=3^2-4\cdot 1\cdot(-3)=21\\
    \lambda_{1,2}&=\frac{-3\pm\sqrt{21}}{2}
  \end{aligned}
\end{equation}
Собственные значения имеют разные знаки, следовательно, точка $(1,1)$
является неустойчивым седлом.

\imageplaceholder{phase_portrait_6.png}{Фазовый портрет системы 6}{fig:phase-system-6}

\subsection{Система 7}

\begin{equation}
  \begin{cases}
    \dot x_1=-x_1^3+x_2^3,\\
    \dot x_2=x_1+3x_3-x_2^3,\\
    \dot x_3=x_1x_3-x_2^3-\sin x_1
  \end{cases}
  \label{eq:system-7}
\end{equation}

Точки равновесия находятся из уравнений
\begin{equation}
  \begin{cases}
    -x_1^3+x_2^3=0,\\
    x_1+3x_3-x_2^3=0,\\
    x_1x_3-x_2^3-\sin x_1=0
  \end{cases}
\end{equation}
Матрица Якоби:
\begin{equation}
  A_7(x_1,x_2,x_3)=
  \begin{pmatrix}
    -3x_1^2 & 3x_2^2 & 0\\
    1 & -3x_2^2 & 3\\
    x_3-\cos x_1 & -3x_2^2 & x_1
  \end{pmatrix}
\end{equation}

Из первого уравнения $x_2=x_1$, а из второго
$x_3=(x_1^3-x_1)/3$. Найдём $x_1$:
\begin{equation}
  x_1^4-3x_1^3-x_1^2-3\sin x_1=0
\end{equation}
Оно имеет два действительных корня: $x_1=0$ и
$x_1\approx3{,}2913$. Соответствующие точки равновесия:
\begin{equation}
  \bm x^{(1)}=(0,0,0),\qquad
  \bm x^{(2)}\approx(3{,}2913;\,3{,}2913;\,10{,}7872)
\end{equation}
В начале координат
\begin{equation}
  \begin{aligned}
    A_7(0,0,0)&=
    \begin{pmatrix}
      0 & 0 & 0\\
      1 & 0 & 3\\
      -1 & 0 & 0
    \end{pmatrix}\\
    \det\bigl(A_7-\lambda I\bigr)
    &=
    \begin{vmatrix}
      -\lambda & 0 & 0\\
      1 & -\lambda & 3\\
      -1 & 0 & -\lambda
    \end{vmatrix}\\
    &=(-\lambda)
    \begin{vmatrix}
      -\lambda & 3\\
      0 & -\lambda
    \end{vmatrix}
    =-\lambda^3=0\\
    \lambda^3&=0\\
    \lambda_1&=\lambda_2=\lambda_3=0
  \end{aligned}
\end{equation}
Метод линеаризации не позволяет определить устойчивость этой точки. Для
второй точки подставим найденные координаты в матрицу Якоби. Получим
\begin{equation}
  \begin{aligned}
    A_7\bigl(\bm x^{(2)}\bigr)
    &\approx
    \begin{pmatrix}
      -32{,}4977 & 32{,}4977 & 0\\
      1 & -32{,}4977 & 3\\
      11{,}7761 & -32{,}4977 & 3{,}2913
    \end{pmatrix}\\
    \det\bigl(A_7-\lambda I\bigr)
    &\approx
    \begin{vmatrix}
      -32{,}4977-\lambda & 32{,}4977 & 0\\
      1 & -32{,}4977-\lambda & 3\\
      11{,}7761 & -32{,}4977 & 3{,}2913-\lambda
    \end{vmatrix}\\
    &=(-32{,}4977-\lambda)
      \left[(-32{,}4977-\lambda)(3{,}2913-\lambda)
      +3\cdot32{,}4977\right]\\
    &\quad
      -32{,}4977\left[(3{,}2913-\lambda)-3\cdot11{,}7761\right]=0\\
    &\Longleftrightarrow
      \lambda^3+61{,}7040\lambda^2
      +907{,}1747\lambda-1348{,}7479=0\\
    \lambda_1&\approx1{,}3585\\
    \lambda_{2,3}&\approx-31{,}5312\pm0{,}5439i
  \end{aligned}
\end{equation}
Поскольку $\lambda_1>0$, вторая точка неустойчива. Тип равновесия -
седло-фокус.

\section{Синтез локально стабилизирующих регуляторов}

Для управляемой системы $\dot{\bm x}=\bm f(\bm x,\bm u)$ необходимо:
\begin{enumerate}[label=\arabic*)]
  \item найти пары $\left(\bm x_{\mathrm{eq}},\bm u_{\mathrm{eq}}\right)$,
  удовлетворяющие $\bm f(\bm x_{\mathrm{eq}},\bm u_{\mathrm{eq}})=0$;
  \item построить линейную модель
  \begin{equation}
    \Delta\dot{\bm x}=A\Delta\bm x+B\Delta\bm u,
    \qquad
    A=\left.\frac{\partial\bm f}{\partial\bm x}\right|_{\mathrm{eq}},
    \quad
    B=\left.\frac{\partial\bm f}{\partial\bm u}\right|_{\mathrm{eq}};
  \end{equation}
  \item проверить управляемость пары $(A,B)$;
  \item выбрать желаемые полюса и найти матрицу обратной связи $K$;
  \item использовать закон управления
  \begin{equation}
    \bm u=\bm u_{\mathrm{eq}}-K(\bm x-\bm x_{\mathrm{eq}})
  \end{equation}
  и проверить устойчивость замкнутой системы численно.
\end{enumerate}

\subsection{Управляемая система 1}

\begin{equation}
  \begin{cases}
    \dot x_1=-x_1+2x_1^3+x_2+\sin u_1,\\
    \dot x_2=-x_1-x_2+3\sin u_2
  \end{cases}
  \label{eq:controlled-system-1}
\end{equation}

Условия равновесия:
\begin{equation}
  \begin{cases}
    -x_1+2x_1^3+x_2+\sin u_1=0,\\
    -x_1-x_2+3\sin u_2=0
  \end{cases}
\end{equation}
Матрицы линеаризованной системы:
\begin{equation}
  A=
  \begin{pmatrix}
    -1+6x_1^2 & 1\\
    -1 & -1
  \end{pmatrix},
  \qquad
  B=
  \begin{pmatrix}
    \cos u_1 & 0\\
    0 & 3\cos u_2
  \end{pmatrix}
\end{equation}

\textbf{Выбранная точка равновесия:}
$\bm x_{\mathrm{eq}}=\ldots$, $\bm u_{\mathrm{eq}}=\ldots$.

\textbf{Желаемые полюса:} $\ldots$.

\textbf{Матрица регулятора:} $K=\ldots$.

\imageplaceholder{control_system_1.png}{Переходные процессы замкнутой управляемой системы 1}{fig:control-system-1}

\subsection{Управляемая система 2}

\begin{equation}
  \begin{cases}
    \dot x_1=x_2+x_1x_2+u^3,\\
    \dot x_2=-x_2+x_2^2-x_1^3+\sin u
  \end{cases}
  \label{eq:controlled-system-2}
\end{equation}

Условия равновесия:
\begin{equation}
  \begin{cases}
    x_2+x_1x_2+u^3=0,\\
    -x_2+x_2^2-x_1^3+\sin u=0
  \end{cases}
\end{equation}
Матрицы линеаризованной системы:
\begin{equation}
  A=
  \begin{pmatrix}
    x_2 & 1+x_1\\
    -3x_1^2 & -1+2x_2
  \end{pmatrix},
  \qquad
  B=
  \begin{pmatrix}
    3u^2\\
    \cos u
  \end{pmatrix}
\end{equation}

\textbf{Выбранная точка равновесия:}
$\bm x_{\mathrm{eq}}=\ldots$, $u_{\mathrm{eq}}=\ldots$.

\textbf{Желаемые полюса:} $\ldots$.

\textbf{Коэффициенты регулятора:} $K=\ldots$.

\imageplaceholder{control_system_2.png}{Переходные процессы замкнутой управляемой системы 2}{fig:control-system-2}


\section{Вывод}

В данном разделе необходимо сопоставить аналитическую классификацию точек
равновесия с фазовыми портретами, сформулировать вывод об устойчивости
системы 4 и оценить качество работы синтезированных регуляторов.
